\section{A classification in every odd prime degree at least seven}
\label{sec:oddprime}

\paragraph{Theorem.}
For each prime $\ell\ge7$, the isomorphism classes of degree-$\ell$ coefficient fields of weight-two CM newforms of trivial nebentypus are in bijection with the imaginary quadratic fields $E$ of class number $\ell$. By the real-embedding argument of Section~\ref{sec:class-three}, this includes every nebentypus. The map assigns to $E$ the real field generated by a root of the Dickson polynomial described below. In particular the number of such fields is $31$ in degree seven, $41$ in degree eleven and $37$ in degree thirteen. These counts use the unconditional class-number enumeration in \href{https://magma.maths.usyd.edu.au/~watkins/papers/smallcn.pdf}{Watkins, final class-number table}.

\paragraph{Restricting the finite character.}
Let $L=EK$ be the Hecke-character value field. Yang's class-number divisibility gives $h_E\mid[L:E]=\ell$ in weight two. If $r$ is the order of the finite component $\eta$, then $r$ is even and $\varphi(r)\mid 2\ell$ because $\mu_r\subset L$. For even $r$ the possibilities are $r=2,4,6$, and $r=2(2\ell+1)$ if $2\ell+1$ is prime. Indeed $\varphi(r)=2\ell$ requires an odd prime factor of $r$; its predecessor divides $2\ell$ and is even. The possibilities for that prime are $3$ and $2\ell+1$. A power of $3$ cannot supply $\ell\ge7$ in the totient, which gives the assertion.

In the last case $q=2\ell+1$ and $L=\mathbb Q(\zeta_q)$, so its unique quadratic subfield is $E=\mathbb Q(\sqrt{-q})$. This field cannot have class number one: none of the prime absolute discriminants $3,7,11,19,43,67,163$ in the class-number-one list has $(q-1)/2$ prime and at least seven. It cannot have class number $\ell$ either. For any negative prime discriminant $-q$ with $q\ge7$, a reduced form $(a,b,c)$ satisfies $a\le\sqrt{q/3}$; there are at most $a$ odd values $-a<b\le a$, by the boundary convention, and $b$ determines $c$. Hence
\[
             h(-q)\le\sum_{1\le a\le\sqrt{q/3}}a
                     <q/3<(q-1)/2.
\]
The same estimate will also exclude $q=\ell$ when $h(-q)=\ell$. If $r=4$ or $6$, then $E$ must be the quadratic cyclotomic field $\mathbb Q(\mu_r)$, since adjoining any other quadratic field to $E$ would give an even subdegree in $L/E$. Both choices have class number one. In class number one, the values on principal ideals already give $L=E(\mu_r)$, a contradiction here. Finally $r=2$ also precludes $h_E=1$, since then $L=E$. Thus $r=2$ and $h_E=\ell$. Genus theory now gives $D_E=-q$ for a single odd prime $q\equiv3\pmod4$ (the single dyadic possibilities have class number one).

\paragraph{Construction, uniqueness, and distinction.}
Choose a nonprincipal degree-one prime ideal $P$ in $E$ away from $q\ell$, with norm $p$, and write $P^\ell=(\alpha)$. Such primes exist by the prime ideal theorem in ideal classes. Define $D_0(X,p)=2$, $D_1(X,p)=X$, $D_n=XD_{n-1}-pD_{n-2}$ for $n\ge2$. A defining polynomial is
\[
                     D_\ell(X,p)-\operatorname{Tr}_{E/\mathbb Q}\alpha.
\]
Indeed, take the Legendre character on $(\mathcal O_E/(\sqrt{-q}))^\times=\mathbb F_q^\times$ as $\eta_0$. Its rational restriction is $\chi_{-q}$ and its sign on $-1$ ensures $\psi((a))=\eta_0(a)a$ is well defined. Extend to the ideal group by choosing $\psi(P)^\ell=\eta_0(\alpha)\alpha$, using that $P$ generates the cyclic class group. Changing the sign of $\psi(P)$ if necessary, write $\beta^\ell=\alpha$. The polynomial $X^\ell-\alpha$ is irreducible over $E$: its reducibility for prime $\ell$ would make $\alpha$ an $\ell$-th power by the binomial irreducibility criterion, forcing $P$ principal. The extension of $\psi$ therefore has value field $E(\beta)$ of degree $\ell$ over $E$. Conversely, for any admissible $\psi$ of finite component order two, use a representative $xP$ prime to its conductor. The value $\psi(xP)/x$ has $\ell$-th power $\pm\alpha$ and generates $L$ over $E$. Thus all resulting value fields are conjugate. Since $\beta\bar\beta=p$, the real element $t=\beta+p/\beta$ satisfies the stated polynomial. If $t$ were rational, $\beta$ would be at most quadratic over $\mathbb Q$ and $\beta^\ell\in E\setminus\mathbb Q$ would force $\beta\in E$. Hence $t$ generates the real coefficient field of prime degree $\ell$.

To see that different $E$ give different coefficient fields, $q$ can be read off from $K$: it is its unique ramified rational prime away from $\ell$. Indeed $M=E(\zeta_\ell,\beta)$ is Galois over $\mathbb Q$ since it contains the radicals above both $\alpha$ and $\bar\alpha$. Away from $q\ell$, its radical extension over $E(\zeta_\ell)$ is unramified: the valuation of $\alpha$ is a multiple of $\ell$ everywhere, and after removing a local $\ell$-th power one has a unit whose $\ell$-th-root polynomial has unit discriminant. At $q\ne\ell$, inertia maps onto the nontrivial inertia of $E/\mathbb Q$, and fixes $\zeta_\ell$. An inertia lift of conjugation sends $\beta$ to $\zeta_\ell^c p/\beta$ for some $c$. On the $\ell$ distinct conjugates
\[
                 t_j=\zeta_\ell^j\beta+\zeta_\ell^{-j}p/\beta
                 \quad(j\bmod\ell)
\]
it acts by $j\mapsto-j-c$, a nonidentity permutation. Hence $q$ ramifies in the splitting field and in $K$ (a normal closure of a field unramified at $q$ is unramified there). This proves injectivity. \hfill$\square$

The class-number divisibility and the general weight-two finiteness theorem were established by \href{https://arxiv.org/abs/math/0301306}{Yang} and \href{https://ems.press/content/serial-article-files/38349}{Gonz\'alez}. The explicit prime-degree bijection above combines those inputs with the realization and ramification arguments. Its proof is independent of a GRH assumption.
